% TOP_PROBLEM: {"id": 1322, "title": "No-Three-in-Line Problem", "category_id": 2, "category": {"id": 2, "name": "combinatorics", "display_name": "Combinatorics", "description": "Counting problems, graph theory, discrete structures.", "slug": "combinatorics", "order_index": 2, "created_at": "2026-07-31T15:26:25.671Z"}, "status": "open"}
% FIRST_POSTED: null
% ENTRY_KIND: statement_only
% Imported from: batch05.tex; UnsolvedMath ID 1322
\documentclass[11pt]{article}
\usepackage[margin=25mm]{geometry}
\usepackage{fontspec}
\setmainfont{Latin Modern Roman}
\usepackage{amsmath,amssymb,mathtools}
\usepackage{unicode-math}
\setmathfont{Latin Modern Math}
\usepackage{xurl,hyperref}
\hypersetup{colorlinks=true,urlcolor=blue}
\setcounter{secnumdepth}{0}
\setlength{\parindent}{0pt}
\setlength{\parskip}{5pt}
\setlength{\emergencystretch}{3em}
\newcommand{\sref}[2]{\hyperlink{source:#1:#2}{[S#2]}}
\newcommand{\cataloguescope}[1]{\paragraph{Recorded target.}#1}
\begin{document}
% UnsolvedMath ID 1322; source catalogue code COMB-004
\setcounter{section}{1321}
\section{No-Three-in-Line Problem}\label{1322}
{\small\sffamily UnsolvedMath ID 1322 · Combinatorics}
\subsection{Definitions and mathematical statement}

\paragraph{Shared notation.}$\mathbb N=\{1,2,\ldots\}$, and $\mathbb Z,\mathbb Q,\mathbb R,\mathbb C$ denote the integers, rationals, reals and complex numbers. Any other conventions are specified locally.


For $n\in\mathbb N$, put $[n]=\{1,\ldots,n\}$ and view $[n]^2\subset\mathbb R^2$ as the square integer grid. Define
\[g(n)=\max\{|A|:A\subset[n]^2,\ \text{no three distinct points of }A\text{ lie on one straight line}\}.\]
Collinearity is ordinary Euclidean collinearity, in every direction, not only horizontal and vertical. Equivalently, for distinct $(a_i,b_i)\in A$, $i=1,2,3$, the determinant of the matrix with rows $(1,a_i,b_i)$ is nonzero.
\paragraph{Statement recorded in the supplied source.}
Determine $g(n)$ exactly for every positive integer $n$, and determine its asymptotic behavior as $n\to\infty$. In particular, is $g(n)=2n$ for every $n\ge2$? \sref{1322}{2}
There can be at most two selected points in each row, so $g(n)\le2n$ for $n\ge2$; the one-point grid has $g(1)=1$. The target is not restricted to a chosen construction or to higher-dimensional grids.
\subsection{Short English statement}
Find the maximum size of a subset of each square integer grid having no three points on any straight line, including its large-grid behavior and the possible $2n$ formula.
\subsection{Sources}
\begin{description}
\item[\hypertarget{source:1322:1}{S1}] ULAM.AI and the UnsolvedMath Contributors, UnsolvedMath: A Curated Collection of Open Mathematics Problems, record 1322 (COMB-004). CC BY 4.0. \url{https://huggingface.co/datasets/ulamai/UnsolvedMath}.
\item[\hypertarget{source:1322:2}{S2}] T. Agama, \emph{On the general no-three-in-line problem}, arXiv:2106.15621v10 (2026), Section 1, p.\ 1, definition of the classical two-dimensional problem. \url{https://arxiv.org/pdf/2106.15621v10}.
\end{description}
\label{end:1322}
\end{document}
